\begin{abstract}

\todo{REWRITE ABSTRACT: Should reflect the new C1-C8 structure. Key changes: (1) Frame around agentic communication being hard, not just "coordination tax". (2) Emphasise the benchmark as a first-class contribution (red-teaming agentic communication). (3) Present solutions (memory, sandboxing, escalation) as proposals validated by experiments, not as finished products. (4) Use "Personal Agent Coordination Layer" framing instead of "AgentOS". (5) Tone down absolute claims.}

Modern knowledge workers spend over 50\% of their time on coordination tasks—scheduling meetings, syncing contexts, and manually managing information flows across organizational boundaries. While AI assistants have made significant strides in productivity, they remain fundamentally isolated, stateless systems incapable of operating securely across trust boundaries. This creates a fundamental bottleneck: effective coordination requires deep personal context, but exposing omniscient agents to external parties creates catastrophic data exfiltration risks.

This thesis presents \textbf{Pulse}, a novel \textbf{Personal Agent Operating System (AgentOS)} that enables secure, stateful AI agents to coordinate autonomously across organizational boundaries. We introduce three core architectural innovations:

\begin{enumerate}
    \item \textbf{Dual-Track Memory}: A memory architecture that strictly separates global self-state from relationship-specific shards, preventing memory contamination across concurrent multi-party interactions.

    \item \textbf{Mountable Context Cells}: A capability-based security mechanism that physically sandboxes agent execution, making prompt injection attacks mathematically ineffective by removing unauthorized API signatures at the system level.

    \item \textbf{Intelligent Escalation Protocol}: An adaptive policy enforcement system that automatically learns social boundaries through relational clustering, reducing manual permission overhead while maintaining zero-trust security guarantees.
\end{enumerate}

Unlike existing single-tenant agent frameworks (OpenClaw, Manus) designed for isolated workspaces, Pulse is architected as a \textbf{network-native system}—analogous to TCP/IP for the agent economy. We demonstrate that through OS-level isolation and relationship-aware memory partitioning, agents can maintain both high collaborative utility and provable security guarantees simultaneously.

We evaluate Pulse through a rigorous benchmark where one host agent interacts with 100 simulated external entities across varying trust tiers. Our system achieves 100\% data exfiltration prevention against jailbreak attacks while maintaining collaborative task completion rates comparable to unrestricted baseline agents. The intelligent escalation protocol reduces manual approval burden by 73\% after the first 20 interactions through precedent learning.

This work establishes the foundational systems architecture for a future where autonomous digital proxies handle cross-boundary coordination, freeing humans to focus exclusively on high-level strategy and creativity.

\end{abstract}